% TODO checklist:
% - [x] Inspected pyproject.toml for ruff, black, mypy configurations
% - [x] Inspected src/logging_setup.py for logging conventions
% - [x] Inspected src/errors/ for error handling patterns
% - [x] Reviewed src/routers/ for API error response patterns
% - [x] Referenced Q&A.md Q46 for technical debt and coding issues
% - [x] Inspected src/utils/ for utility patterns
% - [x] Reviewed README.md for project conventions

\section{Coding Standards and Conventions}
\label{sec:coding-standards}

The Cineca Agentic Platform enforces consistent coding standards across the codebase through automated tooling and established conventions. This section documents the naming conventions, error handling patterns, logging standards, and typing conventions employed throughout the project.

\subsection{Code Formatting and Style}

The project uses automated formatting tools configured in \texttt{pyproject.toml} to ensure consistent code style.

\subsubsection{Black Code Formatter}

Black is configured as the primary code formatter with the following settings:

\begin{lstlisting}[language=Python, caption={Black configuration from \texttt{pyproject.toml}.}]
[tool.black]
line-length = 100
target-version = ["py310", "py311"]
include = "\\.pyi?$"
extend-exclude = """
/(
  docs|examples|ops|db/original-dataset|db/populated
)/
"""
\end{lstlisting}

\subsubsection{Ruff Linting}

Ruff serves as the unified linter and import sorter, replacing flake8, isort, and pylint:

\begin{lstlisting}[language=Python, caption={Ruff lint rules configuration.}]
[tool.ruff]
target-version = "py311"
line-length = 100

[tool.ruff.lint]
select = [
    "E",   # pycodestyle errors
    "F",   # pyflakes
    "W",   # warnings
    "I",   # isort import ordering
    "N",   # pep8-naming
    "UP",  # pyupgrade
    "B",   # flake8-bugbear
    "A",   # flake8-builtins
    "C4",  # flake8-comprehensions
    "SIM", # flake8-simplify
    "PTH", # flake8-use-pathlib
    "RUF", # ruff-specific rules
    "PL",  # pylint-like checks
]
ignore = ["E501"]  # Line length handled by black
\end{lstlisting}

\subsection{Naming Conventions}

The codebase follows Python PEP 8 naming conventions with domain-specific extensions.

\subsubsection{Module and Package Naming}

\begin{table}[ht]
\centering
\caption{Module naming conventions by type.}
\label{tab:module-naming}
\small
\begin{tabular}{lll}
\hline
\textbf{Module Type} & \textbf{Convention} & \textbf{Example} \\
\hline
Router modules & \texttt{snake\_case.py} & \texttt{agent\_runs.py}, \texttt{admin\_db.py} \\
Service modules & \texttt{<domain>\_service.py} or \texttt{<domain>.py} & \texttt{jobs\_service.py} \\
Repository modules & \texttt{<domain>\_repo.py} & \texttt{provider\_repo.py} \\
Model modules & \texttt{snake\_case.py} & \texttt{agent\_run.py} \\
Test modules & \texttt{test\_<module>.py} & \texttt{test\_agents.py} \\
\hline
\end{tabular}
\end{table}

\subsubsection{Class and Function Naming}

\begin{table}[ht]
\centering
\caption{Class and function naming conventions.}
\label{tab:class-naming}
\small
\begin{tabular}{lll}
\hline
\textbf{Element} & \textbf{Convention} & \textbf{Example} \\
\hline
Classes & \texttt{PascalCase} & \texttt{AgentRunRepository} \\
Functions & \texttt{snake\_case} & \texttt{create\_agent\_run()} \\
Private functions & \texttt{\_snake\_case} & \texttt{\_validate\_input()} \\
Constants & \texttt{UPPER\_SNAKE\_CASE} & \texttt{MAX\_STEPS} \\
Type aliases & \texttt{PascalCase} & \texttt{RunResult} \\
FastAPI dependencies & \texttt{get\_<resource>} & \texttt{get\_current\_user()} \\
\hline
\end{tabular}
\end{table}

\subsubsection{MCP Tool Naming}

MCP tools follow a hierarchical naming convention:

\begin{lstlisting}[language=Python, caption={MCP tool naming pattern.}]
# Tool names: <category>.<action>
# Examples:
#   graph.query      - Execute graph queries
#   graph.schema     - Retrieve graph schema
#   security.check   - Validate security constraints
#   system.health    - Component health status

# File location: src/mcp/tools/<category>/<action>.py
# Example: src/mcp/tools/graph/query.py
\end{lstlisting}

\subsection{Error Handling Patterns}

The platform implements structured error handling following the RFC 7807 Problem Details specification.

\subsubsection{Problem Details Response Format}

All API error responses follow the RFC 7807 standard:

\begin{lstlisting}[language=Python, caption={RFC 7807 Problem Details implementation.}]
from pydantic import BaseModel

class ProblemDetail(BaseModel):
    """RFC 7807 Problem Details response."""
    type: str = "about:blank"  # URI reference
    title: str                  # Human-readable summary
    status: int                 # HTTP status code
    detail: str | None = None   # Human-readable explanation
    instance: str | None = None # URI reference to specific occurrence
    
    # Extension fields
    error_code: str | None = None  # Internal error code
    trace_id: str | None = None    # Request trace identifier

# Example error response:
{
    "type": "https://api.cineca.it/errors/validation",
    "title": "Validation Error",
    "status": 422,
    "detail": "Field 'prompt' is required",
    "instance": "/v1/agent-runs",
    "error_code": "VALIDATION_ERROR",
    "trace_id": "abc123..."
}
\end{lstlisting}

\subsubsection{HTTP Exception Handling}

Standard HTTP exceptions are raised and handled consistently:

\begin{lstlisting}[language=Python, caption={HTTP exception raising patterns.}]
from fastapi import HTTPException, status

# Authentication error
raise HTTPException(
    status_code=status.HTTP_401_UNAUTHORIZED,
    detail="Invalid or expired token",
    headers={"WWW-Authenticate": "Bearer"}
)

# Authorization error
raise HTTPException(
    status_code=status.HTTP_403_FORBIDDEN,
    detail=f"Scope '{required_scope}' required"
)

# Not found
raise HTTPException(
    status_code=status.HTTP_404_NOT_FOUND,
    detail=f"Resource '{resource_id}' not found"
)

# Conflict (idempotency)
raise HTTPException(
    status_code=status.HTTP_409_CONFLICT,
    detail="Request with this idempotency key already processed"
)
\end{lstlisting}

\subsubsection{Result Pattern for Service Methods}

Service methods often use a Result pattern for explicit error handling:

\begin{lstlisting}[language=Python, caption={Result pattern for service methods.}]
from dataclasses import dataclass
from typing import Generic, TypeVar

T = TypeVar("T")

@dataclass
class Result(Generic[T]):
    """Result container for service operations."""
    ok: bool
    data: T | None = None
    error: str | None = None
    error_code: str | None = None
    
    @classmethod
    def success(cls, data: T) -> "Result[T]":
        return cls(ok=True, data=data)
    
    @classmethod
    def failure(cls, error: str, error_code: str | None = None) -> "Result[T]":
        return cls(ok=False, error=error, error_code=error_code)

# Usage in services:
async def execute_query(self, cypher: str) -> Result[list[dict]]:
    try:
        results = await self._adapter.execute(cypher)
        return Result.success(results)
    except ValidationError as e:
        return Result.failure(str(e), "VALIDATION_ERROR")
    except TimeoutError:
        return Result.failure("Query timeout", "TIMEOUT")
\end{lstlisting}

\subsection{Logging Conventions}

The platform uses structlog for structured JSON logging with consistent patterns.

\subsubsection{Logger Configuration}

\begin{lstlisting}[language=Python, caption={Structlog configuration from \texttt{src/logging\_setup.py}.}]
import structlog

def setup_logging(level: str = "INFO") -> None:
    """Initialize structlog with JSON output for production."""
    
    structlog.configure(
        processors=[
            structlog.stdlib.filter_by_level,
            structlog.stdlib.add_log_level,
            structlog.stdlib.add_logger_name,
            structlog.processors.StackInfoRenderer(),
            structlog.processors.format_exc_info,
            structlog.stdlib.ProcessorFormatter.wrap_for_formatter,
        ],
        logger_factory=structlog.stdlib.LoggerFactory(),
        cache_logger_on_first_use=True,
        wrapper_class=structlog.stdlib.BoundLogger,
    )
\end{lstlisting}

\subsubsection{Log Event Naming}

Log events follow a hierarchical dot-notation naming convention:

\begin{lstlisting}[language=Python, caption={Log event naming conventions.}]
import structlog

log = structlog.get_logger(__name__)

# Event naming: <component>.<action>.<outcome>
log.info("agent.run.started", run_id=run_id, prompt_length=len(prompt))
log.info("agent.run.completed", run_id=run_id, steps=step_count)
log.warning("agent.run.timeout", run_id=run_id, elapsed_ms=elapsed)
log.error("agent.run.failed", run_id=run_id, error=str(e))

# Component prefixes:
# - agent.*     : Agent orchestration events
# - job.*       : Background job events
# - tool.*      : MCP tool invocation events
# - security.*  : Authentication/authorization events
# - health.*    : Health probe events
# - provider.*  : LLM provider events
\end{lstlisting}

\subsubsection{Context Binding}

Request context is bound to logs for correlation:

\begin{lstlisting}[language=Python, caption={Log context binding pattern.}]
from structlog.contextvars import bind_contextvars, unbind_contextvars

# In middleware:
async def dispatch(self, request: Request, call_next):
    request_id = request.headers.get("x-request-id") or uuid.uuid4().hex
    
    # Bind context for all logs in this request
    bind_contextvars(
        request_id=request_id,
        method=request.method,
        path=request.url.path,
        tenant_id=request.headers.get("x-tenant-id"),
    )
    
    try:
        response = await call_next(request)
        return response
    finally:
        unbind_contextvars("request_id", "method", "path", "tenant_id")
\end{lstlisting}

\subsection{Type Annotations and Mypy}

The codebase uses comprehensive type annotations enforced by mypy:

\begin{lstlisting}[language=Python, caption={Mypy configuration from \texttt{pyproject.toml}.}]
[tool.mypy]
python_version = "3.11"
plugins = ["pydantic.mypy"]
check_untyped_defs = true
disallow_incomplete_defs = true
disallow_untyped_defs = true
no_implicit_optional = true
warn_unused_ignores = true
warn_redundant_casts = true
warn_return_any = true
strict_equality = true
ignore_missing_imports = true

[tool.pydantic-mypy]
init_forbid_extra = true
warn_required_dynamic_aliases = true
init_typed = true
\end{lstlisting}

\subsubsection{Type Annotation Patterns}

\begin{lstlisting}[language=Python, caption={Type annotation examples from the codebase.}]
from typing import Any, TypeVar, Generic
from collections.abc import Callable, Awaitable
from pydantic import BaseModel

# Function signatures with full annotations
async def create_run(
    prompt: str,
    session_id: str | None = None,
    model_id: str | None = None,
    temperature: float = 0.7,
    max_steps: int = 10,
) -> AgentRunResult:
    """Create and execute an agent run."""
    ...

# Generic type variables
T = TypeVar("T", bound=BaseModel)

class Repository(Generic[T]):
    """Generic repository base class."""
    
    def get(self, id: str) -> T | None: ...
    def list(self, filters: dict[str, Any]) -> list[T]: ...
    def create(self, data: T) -> T: ...

# Callback types
ToolHandler = Callable[[dict[str, Any], "ToolContext"], Awaitable[dict[str, Any]]]
\end{lstlisting}

\subsection{API Response Conventions}

API responses follow consistent patterns for pagination, filtering, and metadata.

\subsubsection{List Response Pattern}

\begin{lstlisting}[language=Python, caption={Paginated list response pattern.}]
from pydantic import BaseModel, Field

class PagedResponse(BaseModel, Generic[T]):
    """Standard paginated response wrapper."""
    items: list[T]
    next_page_token: str | None = None
    total_count: int | None = None

# Endpoint example:
@router.get("/agent-runs", response_model=PagedResponse[AgentRunSummary])
async def list_runs(
    page_size: int = Query(default=50, ge=1, le=100),
    page_token: str | None = Query(default=None),
    status: str | None = Query(default=None),
) -> PagedResponse[AgentRunSummary]:
    """List agent runs with pagination."""
    ...
\end{lstlisting}

\subsubsection{Single Resource Response}

\begin{lstlisting}[language=Python, caption={Single resource response with metadata.}]
class ResourceResponse(BaseModel, Generic[T]):
    """Single resource with operational metadata."""
    data: T
    _links: dict[str, str] | None = Field(default=None, alias="links")
    _meta: dict[str, Any] | None = Field(default=None, alias="meta")

# Response headers include:
# - ETag: "abc123" for caching
# - X-Request-ID: correlation identifier
# - X-Trace-Id: distributed trace ID
\end{lstlisting}

\subsection{Configuration Conventions}

Configuration follows the Twelve-Factor App methodology with environment variables:

\begin{lstlisting}[language=Python, caption={Configuration conventions from \texttt{src/config.py}.}]
from pydantic_settings import BaseSettings, SettingsConfigDict
from pydantic import Field

class Settings(BaseSettings):
    """Application configuration from environment variables."""
    
    # Naming: UPPER_SNAKE_CASE matching env vars
    APP_ENV: str = Field(default="dev")
    APP_HOST: str = Field(default="0.0.0.0")
    APP_PORT: int = Field(default=8000)
    
    # Grouped by concern with prefixes:
    # DB_* for database settings
    DB_HOST: str = Field(default="postgres")
    DB_PORT: int = Field(default=5432)
    
    # REDIS_* for cache settings
    REDIS_URL: str = Field(default="redis://redis:6379/0")
    
    # OIDC_* for authentication
    OIDC_ISSUER: str | None = Field(default=None)
    
    # Settings configuration
    model_config = SettingsConfigDict(
        env_file=".env",
        env_file_encoding="utf-8",
        extra="allow",
    )
\end{lstlisting}

\subsection{Documentation Conventions}

Code documentation follows Google-style docstrings:

\begin{lstlisting}[language=Python, caption={Docstring conventions.}]
async def execute_run(
    self,
    prompt: str,
    options: RunOptions | None = None,
) -> AgentRunResult:
    """Execute an agent run with the given prompt.
    
    This method orchestrates the complete agent execution lifecycle,
    including intent classification, step planning, and execution.
    
    Args:
        prompt: The user's natural language input.
        options: Optional configuration for the run.
    
    Returns:
        AgentRunResult containing the final output and execution metrics.
    
    Raises:
        ValidationError: If the prompt is empty or invalid.
        TimeoutError: If execution exceeds the configured timeout.
        AuthorizationError: If the user lacks required permissions.
    
    Example:
        >>> result = await orchestrator.execute_run(
        ...     "What users work at CINECA?",
        ...     RunOptions(max_steps=5)
        ... )
        >>> print(result.output)
    """
\end{lstlisting}

\subsection{Security Code Conventions}

Security-sensitive code follows additional conventions:

\begin{lstlisting}[language=Python, caption={Security coding conventions.}]
# 1. Never log sensitive data
log.info("auth.token.validated", user_id=user_id)  # OK
log.info("auth.token", token=token)  # BAD - logs token

# 2. Use constant-time comparisons for secrets
import hmac
def verify_signature(expected: bytes, actual: bytes) -> bool:
    return hmac.compare_digest(expected, actual)

# 3. Sanitize user input before logging
def safe_log_input(user_input: str) -> str:
    return user_input[:100].replace("\n", " ")

# 4. Use parameterized queries (never string interpolation)
# Good:
query = "MATCH (u:User {id: $user_id}) RETURN u"
params = {"user_id": user_id}

# Bad:
query = f"MATCH (u:User {{id: '{user_id}'}}) RETURN u"
\end{lstlisting}

